\begin{statement}{A.5.17}{Proved}{Isolated singularities are excluded and suspension germs are transverse products}
Let $f:Y\to C_0$ have the hypotheses of A.5.15.
No local equation $f=u h^m$, with $u$ a holomorphic
unit and $m\ge1$, can have a reduced equation $h$
with an isolated critical point at the origin.
Suppose its local equation at a point of a fibre is
\[
             f=(x^a-g(y,z))^m,
             \qquad a\ge2,\quad m\ge1,             \tag{A.J.38}
\]
where $g\ne0$ and $g(0,0)=0$. Then the reduced
plane germ $\{g=0\}_{\rm red}$ is smooth. If
$b$ is its multiplicity in $\operatorname{div}(g)$,
there are local holomorphic coordinates in which
$f=(x^a-y^b)^m$.
Consequently, at every point where the full fibre
is a uniform multiple of its reduced support and
the reduced equation has a nonzero Hessian, its
critical surface germ is a constant
transverse $A_{b-1}$ product. A nonordinary change
of transverse type at a point with nonzero Hessian
is excluded.
\end{statement}
\begin{proof}
We first prove the isolated assertion, including its
comparison with the actual global fibre. Replace $h$
by $H=u^{1/m}h$, taking a holomorphic root on a small
coordinate neighbourhood. Then $f=H^m$ exactly.
This replacement preserves an isolated critical point
as a germ. Indeed every local irreducible component
of the critical locus of $H$ has constant $H$-value,
necessarily zero when it contains the origin. On
$H=0$ one has $dH=u^{1/m}dh$. Thus its critical
locus has the same zero-dimensional germ as that
of $h$.

Choose a small closed ball $K\subset\C^3$ such that
$H^{-1}(0)$ is smooth along its boundary intersection
and transverse to $\partial K$.
The standard isolated-singularity Milnor-ball
construction supplies this choice; see
\cite[Lemma 6.1 and Appendix B]{Milnor1968}.
The restriction $H|_{\partial K}$ is a submersion
at its zero set. Compactness therefore supplies
$\eta>0$ for which it is a submersion wherever
$|H|<2\eta$ on the boundary. For the linear
deformation
\[
 H_\lambda(w)=H(w)-\lambda_1w_1
                         -\lambda_2w_2-\lambda_3w_3,
                                                        \tag{A.J.38a}
\]
this boundary submersion persists for all parameters
in a sufficiently small polydisc $P$ and all values
of modulus less than $\eta$. The gradient map of
$H$ has an isolated zero and consequently is a finite
holomorphic map germ of positive local degree. Choose
a small regular value $\lambda$ of that map. Its
nonempty finite inverse image consists of critical
points of $H_\lambda$ with invertible Hessian. All
these points and their critical values tend to the
origin as $\lambda\to0$.

Remove the critical-value discriminant from
$P\times\Delta_\eta$. The relative critical locus
over this domain misses the boundary, so its
projection is proper and its image is a proper
complex analytic subset. Its complement is path
connected. The projection of the corresponding
smooth level sets in $K$ is a proper submersion
with boundary, and hence an Ehresmann bundle.
At any one of the nondegenerate critical points,
the holomorphic Morse lemma gives
$H_\lambda-H_\lambda(p)=X^2+Y^2+Z^2$.
A local Milnor fibre at a sufficiently small regular
value embeds in the corresponding level set in $K$.
The bundle transports it, by an oriented embedding,
into $K\cap H^{-1}(\epsilon_0)$ for a nonzero small
$\epsilon_0$, which is a submanifold of the actual
global torus $Y_{\epsilon_0^m}$.

For a positive real local Milnor value, the real
locus $X,Y,Z\in\mathbb R$ is a two-sphere. It is
Lagrangian for the standard K\"ahler form; multiplication
by $\mathrm i$ identifies its tangent bundle with
its normal bundle with the opposite oriented Euler
class. Its self-intersection is therefore
$-\chi(S^2)=-2$. Equivalently this is the case
$a=2$, $b=c=1$ of the explicit cycle in A.5.15,
after the linear change $Y^2+Z^2=(Y+\mathrm iZ)(Y-\mathrm iZ)$.
The transported sphere has the same square in the
actual torus, whose $\pi_2$ is zero. This is the
required contradiction. No suspension representation
of the original isolated singularity has been used.

We prove the comparison of Milnor fibres needed for
the argument. Factor the analytic plane function as
\[
                   g=u\prod_{i=1}^r h_i^{b_i},
\]
where the $h_i$ are distinct reduced irreducible
germs and $u$ is a unit. Deform each $h_i$
holomorphically to $h_{i,s}$ in a finite-dimensional
parameter polydisc $P$, and retain its exponent:
\[
 g_s=u\prod_i h_{i,s}^{b_i},\qquad
 H_s=x^a-g_s.                                    \tag{A.J.39}
\]
We use only sufficiently small parameters.

Choose a small closed ball $K\subset\C^2_{y,z}$
such that each $h_i=0$ meets $\partial K$
transversely in its smooth part, and their boundary
links are disjoint. Choose its representative so
that every critical value of $g$ on $K$ is zero.
This is possible because the critical locus of a
holomorphic function has finitely many local
irreducible components, the function is constant
on each of them, and each component of the germ
contains the origin. The boundary smoothness,
transversality and disjointness persist for
$h_{i,s}$ by compactness and the implicit function
theorem. Near any of those links, the other factors
are units, so there are parameter-dependent local
coordinates $(v,w)$ with $g_s=v^{b_i}$, obtained
by taking a root of that unit. The boundary
$\partial K$ is transverse to the $w$ direction.
Consequently the restriction of $g_s$ to
$\partial K$ is a submersion at every point where
$0<|g_s|<\eta_0$, for a uniform sufficiently small
$\eta_0>0$. Away from these link neighbourhoods
$|g_s|$ has a uniform positive lower bound on the
boundary. This proves the asserted uniform
transversality, including at parameters for which
the germs inside $K$ change their singularities.

Take a small disc $D_x$ of radius $\rho$ with
$\rho^a<\eta_0/2$, then shrink $P$ so that every
critical value of $g_s$ on $K$ has modulus less than
$\rho^a/4$. This shrinking follows by compactness:
a sequence of critical points with parameters
tending to zero would otherwise limit to a
nonzero critical value of $g$ in the chosen small
representative. The latter has no such value.
Fix $0<\eta<\rho^a/4$ and put
$W=\overline D_x\times K$, a compact manifold
with corners contained in the original coordinate
chart in $Y$.

Every $H_s^{-1}(\epsilon)$ with
$0<|\epsilon|<\eta$ is transverse to this boundary.
On the face $\partial D_x\times K$, the equation
makes $|g_s|>3\rho^a/4$, so $dg_s\ne0$; its plane
directions give the submersion. On the face
$\overline D_x\times\partial K$, the $x$
derivative gives it when $x\ne0$. At $x=0$ we
have $g_s=-\epsilon\ne0$, and the preceding
boundary calculation gives it. At the corner
$|x|=\rho$, that same calculation for $g_s$
supplies both real target directions. The corner
stratum is therefore also transverse. We use
Ehresmann's theorem with corners: local lifts of
base tangent vectors can be chosen tangent to
each boundary face by these submersion statements,
then patched by a partition of unity and
integrated on the compact fibres. Thus no change
of the chosen domain $W$ is required.

Inside $P\times\{0<|\epsilon|<\eta\}$ remove the
critical-value discriminant $\mathcal D$ of the
maps $H_s$ on $W$. This is a proper complex
analytic subset. The relative critical locus has
no point on the boundary at such a value, so its
projection is proper over every compact subset
of this parameter domain, and the proper mapping
theorem applies. For fixed $s$ its critical
values are discrete, since a holomorphic function
is constant on each connected critical component.
The complement of a proper complex analytic
subset of a connected complex manifold is path
connected. The transversality just proved and
the implicit function theorem give a proper
smooth submersion, including its boundary,
\[
 \{(s,p):p\in W,\ (s,H_s(p))\in
                    (P\times\Delta_\eta^*)\setminus\mathcal D\}
 \longrightarrow
       (P\times\Delta_\eta^*)\setminus\mathcal D,
 \quad (s,p)\longmapsto(s,H_s(p)).                 \tag{A.J.40}
\]
Ehresmann's theorem identifies all its compact
fibres by orientation-preserving diffeomorphisms
of manifolds with corners.
In particular a sufficiently small local Milnor
fibre of an interior singularity of some $H_s$,
at a value near zero, embeds into a smooth fibre
in (A.J.40), hence into
$W\cap H_0^{-1}(\epsilon_0)$ for a sufficiently
small nonzero $\epsilon_0$. This last space is
an actual submanifold of the global torus
$Y_{\epsilon_0^m}$ by (A.J.38); no identification
of it with a standard Milnor fibre is needed.
Choose the value to avoid every other
critical value, and choose its local Milnor ball
in the interior of $W$. This is the precise
comparison we use. Arbitrary deformations which
change the exponents along the boundary links
are not used in (A.J.40).

Suppose first that $r\ge2$. Choose two distinct
factors $h_i,h_j$. Their map
$(h_i,h_j):(\C^2,0)\to(\C^2,0)$ is finite,
since their common zero germ is zero dimensional.
Its positive local degree is their intersection
multiplicity. A sufficiently small regular value
$(s_i,s_j)$ consequently has nonempty finite
inverse image. At each inverse-image point,
$h_i-s_i$ and $h_j-s_j$ are transverse coordinates.
Choose the remaining small shifts to avoid those
finitely many points. The deformation (A.J.39)
then has an interior point at which
\[
            H_s=x^a-U v^{b_i}w^{b_j},
\]
with $U$ a unit independent of $x$. Replacing
$v$ by $vU^{1/b_i}$ makes this exactly
$x^a-v^{b_i}w^{b_j}$ by an invertible holomorphic
coordinate change. A.5.15 constructs a spherical
class of nonzero square in its local Milnor fibre.
The embedding (A.J.40) puts that class in the
actual smooth torus fibre of $f=H_0^m$.
Preservation of intersection
numbers and $\pi_2$-vanishing give the same
contradiction as in A.5.15. Thus $r=1$.

Suppose its reduced factor $h=h_1$ is singular.
Write $g=u h^b$ and absorb a local $b$th root of
$u$ into $h$. This leaves its reduced zero germ
unchanged. Since $h$ is reduced, its gradient
has an isolated zero at the origin: a critical
curve would have constant value zero and would
be a repeated component of $h$. The finite
gradient map has a positive local degree.
Choose a small regular value
$(\lambda_1,\lambda_2)$ of this map. Then
\[
        h_s=h-\lambda_1y-\lambda_2z-\lambda_0
\]
has nondegenerate critical points near zero;
regularity of the gradient value says exactly
that their Hessians are invertible. Choose
$\lambda_0$ so that their critical values are
nonzero. At one such point $p$ put
$v=h_s(p)\ne0$. The Hessian of $h_s^b$ is
$b v^{b-1}$ times the Hessian of $h_s$, and
is invertible. The holomorphic Morse lemma on
the two plane variables identifies the germ of
$H_s-H_s(p)$ with
\[
                         x^a-Y^2-Z^2.
\]
A linear change of $(Y,Z)$ gives the Hirzebruch
germ $x^a-YZ$, already excluded by A.5.15.
Its Milnor fibre embeds by (A.J.40), producing
the same contradiction. Hence $h$ is smooth.
Taking $y=h u^{1/b}$ as a coordinate proves the
asserted transverse product.

For the Hessian consequence, take a local $m$th
root of the actual map, absorbing its unit.
Its reduced equation has the same nonzero
Hessian rank as every other reduced equation
at the critical point. The holomorphic splitting
lemma splits off a nondegenerate square and
gives $f=(x^2-g(y,z))^m$ exactly. If $g=0$,
the equation would not be reduced, so $g\ne0$.
At a critical point $dg(0)=0$. The result just
proved gives $g=y^b$ with $b\ge2$ in suitable
holomorphic coordinates, which is the constant
transverse $A_{b-1}$ product.
\end{proof}


% A.5.18 from part_A6_tacnode_bounds.tex

